\origpage{156}{ROBERTSON AND SEYMOUR}
\noindent has order $\ge 2$. Now $(K, G \setminus v)$ has order 0, and $(G \setminus v, K) \notin \mathscr{T}$ by the second axiom, since $(G \setminus v) \cup K = G$. Thus $(K, G \setminus v) \in \mathscr{T}$. But $(K \cup A_1, (G \setminus v) \cap B_1)$ has order at most the order of $(A_1, B_1)$ and hence is in $\mathscr{T}$ by (2.2), contrary to the maximality of $|V(A_1)|$, as required. $\square$

For an edge $e$ of a hypergraph $G$, the \textit{ends} of $e$ are the vertices of $G$ incident with $e$, and the \textit{size} of $e$ is the number of ends of $e$.

\stmt
(2.4) \textit{Let $\theta \ge 1$, and let $e$ be an edge of $G$ with size $\ge \theta$. Let $\mathscr{T}$ be the set of all separations $(A, B)$ of $G$ of order $< \theta$ with $e \in E(B)$. Then $\mathscr{T}$ is a tangle of order $\theta$.}

\textit{Proof.} The first two axioms are clear. For the third, let $(A, B) \in \mathscr{T}$. Then $V(A \cap B)$ does not contain every end of $e$ since $|V(A \cap B)| < \theta$, and yet $e \in E(B)$, and so $V(A) \ne V(G)$. This completes the proof. $\square$

We remark

\stmt
(2.5) \textit{$G$ has a tangle if and only if $V(G) \ne \emptyset$.}

\textit{Proof.} If $v \in V(G)$, let $\mathscr{T}$ be the set of all separations $(A, B)$ of $G$ of order 0 with $v \in V(B)$. Then $\mathscr{T}$ is a tangle of order 1, as is easily seen. Conversely, since every tangle has order $\le |V(G)|$, if $G$ has a tangle then $V(G) \ne \emptyset$. $\square$

For graphs, we can extend (2.5) as follows.

\stmt
(2.6) \textit{If $G$ is a graph, the tangles in $G$ of order 1 are in 1--1 correspondence with the connected components of $G$, and those of order 2 are in 1-1 correspondence with the blocks of $G$ which have a non-loop edge.}

(A block of a graph is a maximal connected subgraph any two distinct edges of which are in a circuit.)

\textit{Proof.} Since we do not need the result, we merely sketch the proof. Any $v \in V(G)$ yields a tangle of order 1 as in (2.5), and it is easy to see that every tangle of order 1 arises this way, and distinct $v, v' \in V(G)$ yield the same tangle if and only if $v$ and $v'$ are in the same component of $G$. For order 2, any non-loop edge yields a tangle of order 2, by (2.4), and again, it is easy to see that every order 2 tangle arises this way, and two edges yield the same tangle if and only if they are in the same block. \fix{We give the details. Besides (2.2)--(2.4) we use (2.7), (2.9) and the extreme separations defined before (2.8), whose proofs do not use (2.6). Let $K$ be the null graph. For every tangle $\mathscr{T}$ in $G$, the pair $(K, G)$ is in $\mathscr{T}$, since $(K, G)$ and $(G, K)$ are separations of order 0 and $(G, K) \notin \mathscr{T}$ by the third axiom.

For $v \in V(G)$ let $\mathscr{T}_v$ be the set of all separations $(A, B)$ of $G$ of order 0 with $v \in V(B)$. This is the tangle of order 1 of (2.5). Indeed, if $(A, B)$ has order 0 then $v$ is in exactly one of $V(A)$, $V(B)$, which gives the first axiom, and if $(A_i, B_i) \in \mathscr{T}_v$ $(1 \le i \le 3)$ then $v$ is in none of $V(A_1)$, $V(A_2)$, $V(A_3)$, which gives the second and third.

\hspace*{12pt}(1) \textit{Every tangle $\mathscr{T}$ in $G$ of order $1$ equals $\mathscr{T}_v$ for some $v \in V(G)$.}

Choose $(A, B) \in \mathscr{T}$ with $|V(A)| + |E(A)| - |V(B)| - |E(B)|$ maximum (this is possible since $(K, G) \in \mathscr{T}$). Then $(A, B)$ is extreme, since for $(A', B') \in \mathscr{T}$ with $A \subseteq A'$ and $B' \subseteq B$ the maximality gives $A' = A$ and $B' = B$. By the third axiom there exists $v \in V(G) - V(A)$, and $v \in V(B)$ since $A \cup B = G$. Let $(A', B') \in \mathscr{T}$. A vertex in both $A \cup A'$ and $B \cap B'$ is in $A \cap B$ or in $A' \cap B'$, and so $(A \cup A', B \cap B')$ has order 0. Hence it is in $\mathscr{T}$ by (2.2), and $B \cap B' = B$ since $(A, B)$ is extreme. Thus $v \in V(B')$, and so $\mathscr{T} \subseteq \mathscr{T}_v$. Conversely, let $(A', B') \in \mathscr{T}_v$. If $(B', A') \in \mathscr{T}$ then $v \in V(A')$ by what we just proved, and so $v \in V(A' \cap B')$, which is impossible since $(A', B')$ has order 0. Hence $(A', B') \in \mathscr{T}$ by the first axiom. This proves (1).

\hspace*{12pt}(2) \textit{For $u, v \in V(G)$, $\mathscr{T}_u = \mathscr{T}_v$ if and only if $u$ and $v$ are in the same component of $G$.}

Let $(A, B)$ be a separation of order 0. An edge of $G$ with an end in $V(B)$ is not in $E(A)$, since all ends of an edge of $A$ are in $V(A)$ and no vertex is in both. So it is in $E(B)$, and all its ends are in $V(B)$. Hence if $u$ and $v$ are joined by a path of $G$, then $u \in V(B)$ if and only if $v \in V(B)$, and so $\mathscr{T}_u = \mathscr{T}_v$. Now let $u$ and $v$ be in different components, let $A$ be the component that contains $u$, and let $B$ be the subgraph formed by all the other vertices and edges of $G$. Then $(A, B)$ is a separation of order 0 since no edge of $B$ has an end in $V(A)$, and $(A, B) \in \mathscr{T}_v - \mathscr{T}_u$. This proves (2).

By (1) and (2), mapping each component $C$ of $G$ to $\mathscr{T}_v$, where $v \in V(C)$, gives a 1-1 correspondence between the components of $G$ and the tangles of order 1. (By (2) the image does not depend on the choice of $v$ and distinct components have distinct images, and by (1) every tangle of order 1 is an image.)

For a non-loop edge $e$ let $\mathscr{T}_e$ be the set of all separations $(A, B)$ of $G$ of order $\le 1$ with $e \in E(B)$. Since $e$ has size 2, $\mathscr{T}_e$ is a tangle of order 2 by (2.4).

\hspace*{12pt}(3) \textit{Every tangle $\mathscr{T}$ in $G$ of order $2$ equals $\mathscr{T}_e$ for some non-loop edge $e$.}

For $X \subseteq E(G)$, let $L(X)$ be the subgraph with edge set $X$ whose vertices are the ends of the edges in $X$ and the vertices incident with no edge, and let $R(X)$ be the subgraph with edge set $E(G) - X$ whose vertices are the ends of the edges not in $X$. Then $(L(X), R(X))$ is a separation of $G$. If $(A, B) \in \mathscr{T}$ then $(L(E(A)), R(E(A))) \in \mathscr{T}$. For a vertex in both $L(E(A))$ and $R(E(A))$ is an end of an edge of $A$ and of an edge of $B$, and so it is in $V(A \cap B)$. Hence $(L(E(A)), R(E(A)))$ has order $\le 1$, and it is in $\mathscr{T}$ by (2.9)(iii). In particular $(L(\emptyset), R(\emptyset)) \in \mathscr{T}$, since $(K, G) \in \mathscr{T}$.

We claim that some edge $e$ of $G$ is in $E(B)$ for all $(A, B) \in \mathscr{T}$. Suppose not. Then every edge is in $E(A)$ for some $(A, B) \in \mathscr{T}$, and so it is in some $X \subseteq E(G)$ with $(L(X), R(X)) \in \mathscr{T}$. Choose a set $\mathscr{S}$ of subsets $X$ of $E(G)$ with union $E(G)$, such that $(L(X), R(X)) \in \mathscr{T}$ for all $X \in \mathscr{S}$, with $|\mathscr{S}|$ minimum. If $|\mathscr{S}| \le 3$, then the members of $\mathscr{S}$ and copies of $\emptyset$ give three members of $\mathscr{T}$ the union of whose first terms contains every edge of $G$, contrary to (2.3). So $|\mathscr{S}| \ge 4$. Choose $X_1 \in \mathscr{S}$, and put $(A_1, B_1) = (L(X_1), R(X_1))$. Let $X' \in \mathscr{S} - \{X_1\}$, and put $(A', B') = (L(X'), R(X'))$. Then $(A_1 \cup A', B_1 \cap B')$ has order $\ge 2$. For otherwise it is in $\mathscr{T}$ by (2.2), and so $(L(X_1 \cup X'), R(X_1 \cup X')) \in \mathscr{T}$ by the previous paragraph, since $E(A_1 \cup A') = X_1 \cup X'$. But then replacing $X_1$ and $X'$ by $X_1 \cup X'$ in $\mathscr{S}$ gives a smaller set, contrary to the choice of $\mathscr{S}$. Now every vertex of $(A_1 \cup A') \cap (B_1 \cap B')$ is either in $V(A_1 \cap B_1) \cap V(B')$ or in $V(A' \cap B') - V(A_1)$. Each of these two sets has at most one member, since both separations have order $\le 1$, and together they have at least two. So each has exactly one. Hence $V(A_1 \cap B_1) = \{u\}$ for some $u \in V(B')$, and the only vertex of $A' \cap B'$ is not in $V(A_1)$, so it is not $u$. Since $u \in V(B')$, it follows that $u \notin V(A')$. Here $u$ does not depend on $X'$, since $V(A_1 \cap B_1) = \{u\}$. Now $u \in V(B_1)$, and so $u$ is an end of some edge $f \notin X_1$. Since the union of $\mathscr{S}$ is $E(G)$, $f \in X'$ for some $X' \in \mathscr{S} - \{X_1\}$, and then $u \in V(L(X'))$, a contradiction. This proves our claim.

Choose $e$ as in the claim. If $e$ is a loop, then $e$ has size 1, and so $(K_e, G \setminus e) \in \mathscr{T}$ by (2.7), where $K_e$ is formed by $e$ and its end. This is impossible since $e \notin E(G \setminus e)$. So $e$ is not a loop, and $\mathscr{T} \subseteq \mathscr{T}_e$ by the claim. Conversely, let $(A, B) \in \mathscr{T}_e$. Then $e \notin E(A)$, and so $(B, A) \notin \mathscr{T}$ by the claim. Hence $(A, B) \in \mathscr{T}$ by the first axiom. This proves (3).

\hspace*{12pt}(4) \textit{If two distinct edges $e_1, e_2$ of $G$ are in a circuit $C$, then $\mathscr{T}_{e_1} = \mathscr{T}_{e_2}$.}

Since $C$ has two edges, it is not a loop, and so every edge of $C$ has two ends and every vertex of $C$ is incident with exactly two edges of $C$. Let $(A, B)$ be a separation of order $\le 1$. For $v \in V(C)$ let $d(v)$ be the number of edges of $C$ in $E(A)$ incident with $v$. The sum of $d(v)$ over all $v \in V(C)$ is twice the number of edges of $C$ in $E(A)$, and so it is even. If $v \in V(A) - V(B)$ then $d(v) = 2$, since an edge of $B$ at $v$ would put $v$ in $V(B)$, and if $v \notin V(A)$ then $d(v) = 0$. Since at most one vertex of $C$ is in $V(A \cap B)$, it follows that $d(v)$ is even for that vertex too. So at each vertex of $C$, either both edges of $C$ are in $E(A)$ or neither is. Since $C$ is connected, either every edge of $C$ is in $E(A)$ or none is. Since every edge of $G$ is in exactly one of $E(A)$, $E(B)$, it follows that $e_1 \in E(B)$ if and only if $e_2 \in E(B)$. This proves (4).

\hspace*{12pt}(5) \textit{If $e_1, e_2$ are distinct non-loop edges with $\mathscr{T}_{e_1} = \mathscr{T}_{e_2}$, then some circuit of $G$ contains both.}

Let $Y_1$, $Y_2$ be the sets of ends of $e_1$, $e_2$, and let $G' = G \setminus \{e_1, e_2\}$. Suppose that $G'$ has a separation $(G_1, G_2)$ of order $\le 1$ with $Y_1 \subseteq V(G_1)$ and $Y_2 \subseteq V(G_2)$. Adding $e_1$ to $G_1$ and $e_2$ to $G_2$ gives a separation $(A, B)$ of $G$ of order $\le 1$ with $e_1 \in E(A)$ and $e_2 \in E(B)$. Then $(A, B) \in \mathscr{T}_{e_2} = \mathscr{T}_{e_1}$, and so $e_1 \in E(B)$, a contradiction. Hence by Menger's theorem there are two disjoint paths of $G'$ from $Y_1$ to $Y_2$. Since they are disjoint and $|Y_1| = |Y_2| = 2$, one of them joins one end of $e_1$ to one end of $e_2$, the other joins the other two ends, and neither has an internal vertex in $Y_1 \cup Y_2$. So the two paths together with $e_1$ and $e_2$ form a circuit of $G$. This proves (5).

For a non-loop edge $e$, let $H_e$ be the subgraph of $G$ whose edges are the non-loop edges $f$ with $\mathscr{T}_f = \mathscr{T}_e$, and whose vertices are the ends of these edges.

\hspace*{12pt}(6) \textit{For every non-loop edge $e$, $H_e$ is a block of $G$, and every block of $G$ that contains $e$ equals $H_e$.}

Let $f_1, f_2$ be distinct edges of $H_e$. Then $\mathscr{T}_{f_1} = \mathscr{T}_{f_2}$, and so some circuit $C$ of $G$ contains $f_1$ and $f_2$, by (5). Every edge of $C$ is a non-loop edge with the same tangle as $f_1$, by (4), and so $C \subseteq H_e$. Hence any two distinct edges of $H_e$ are in a circuit of $H_e$. Also $H_e$ is connected, since any two of its vertices are ends of edges $f_1, f_2$ of $H_e$, and either $f_1 = f_2$ or $f_1, f_2$ are in a circuit of $H_e$. Now let $H$ be a connected subgraph of $G$ that contains $e$, any two distinct edges of which are in a circuit of $H$. We show that $H \subseteq H_e$. Each edge $f \ne e$ of $H$ is in a circuit of $H$ together with $e$. This circuit is not a loop, and so $f$ is a non-loop edge with $\mathscr{T}_f = \mathscr{T}_e$, by (4). So $E(H) \subseteq E(H_e)$. Each vertex of $H$ is incident with an edge of $H$, since $H$ is connected and contains the edge $e$. So $V(H) \subseteq V(H_e)$, and $H \subseteq H_e$. Applying this to a subgraph $H$ that contains $H_e$ shows that $H_e$ is maximal, and so $H_e$ is a block. Applying it to a block $H$ that contains $e$ gives $H \subseteq H_e$, and so $H = H_e$ by the maximality of $H$. This proves (6).

For non-loop edges $e, e'$, $H_e = H_{e'}$ if and only if $\mathscr{T}_e = \mathscr{T}_{e'}$, by the definition of $H_e$ and since $e \in E(H_e)$. Map each block $H$ with a non-loop edge to $\mathscr{T}_e$, where $e$ is a non-loop edge of $H$. By (6), $H = H_e$ for every such $e$, and so the image does not depend on the choice of $e$, and distinct blocks have distinct images. By (3) and (6), every tangle of order 2 is $\mathscr{T}_e$ for some non-loop edge $e$, and so it is the image of the block $H_e$. This gives the 1-1 correspondence for order 2.}\fixnote{fix:2.6-proof-added}{The paper only sketches this proof. The details above are the argument used in the formalization.} $\square$

One might speculate that in a graph, the tangles of order $d$ correspond to the long-sought ``$d$-connected components,'' but that possibility is not further explored here.
